\chapter{Introduction}

\section{Context and motivation}

Fair exchange is the problem of allowing two parties to exchange valuable assets in such a
way that no honest participant is left at a disadvantage. In the setting considered in this
project, a vendor \(V\) owns a digital item \(x\), a buyer \(B\) wants to buy it for a price
\(y\), and both parties know in advance a public description predicate \(\mathsf{desc}\) and
an expected digest \(d\). The intended outcome is simple to state: either \(B\) obtains an
item \(x'\) such that \(\mathsf{desc}(x') = d\) and \(V\) receives the payment, or the exchange
is cancelled without giving either side an unfair advantage.

This problem is classical in cryptography, but it remains difficult to deploy in practice.
Without an external trusted party, perfect two-party fair exchange is impossible in the
general case. Smart contracts provide a natural alternative: the blockchain can act as a
public, deterministic, and economically enforceable adjudicator. However, using a blockchain
introduces another difficulty. Every on-chain action costs gas, and these costs can make the
exchange unattractive, especially when the exchanged item is small, when one participant has
no cryptocurrency available, or when a dispute forces expensive verification on-chain.

The Sponsored Fair Exchange protocol, denoted \sox{} in this report, addresses this tension
by combining three ideas. First, it follows an optimistic design: when both parties behave
honestly, the exchange should finish with very little on-chain computation. Second, it keeps
a dispute mechanism able to resolve disagreements by verifying only a logarithmic number of
positions in a computation trace. Third, it introduces sponsors, whose role is to cover gas
costs under predefined incentives. The goal is not only fairness, but also economic
accessibility: the vendor and the buyer should not need to manage gas in the common case.

The starting point of this semester project was an existing implementation of \sox{},
documented in the previous implementation report~\cite{hana2025} and based on the
protocol paper~\cite{sox2026}. That implementation already represented a substantial
proof of concept. It contained a Rust/WebAssembly off-chain pipeline, Solidity contracts for
the optimistic and dispute phases, an Electron/Next.js interface, and an ERC-4337-inspired
transaction sponsorship path. It also introduced a compact block-based circuit representation
for large files, with 64-byte gates and Merkle commitments over the circuit, ciphertext, and
intermediate values.

The purpose of the present project is different from reimplementing \sox{} from scratch. The
objective is to take this initial system seriously as a research prototype, analyze where its
implementation diverges from the most gas-efficient protocol variants, and modify it so that
the main variants suggested by the protocol analysis can be executed and measured. In other
words, the project is both an implementation project and an experimental methodology project:
the code must work, but the measurements must also correspond to well-defined protocol
scopes.

\section{Terminology used throughout the report}

A recurring source of confusion in this project is that several implementations coexist in
the repository and they do not measure the same thing. To avoid ambiguity, this report uses
the following terminology consistently.

\begin{table}[htbp]
  \centering
  \begin{tabularx}{\linewidth}{>{\raggedright\arraybackslash}p{0.28\linewidth}>{\raggedright\arraybackslash}X}
    \toprule
    Term & Meaning in this report \\
    \midrule
    Initial implementation &
    The implementation inherited from the previous semester project. It is the baseline used
    to understand the existing architecture and to reproduce the original costs. \\
    Monolithic implementation &
    The intermediate implementation in which several variants are supported by the same
    large contracts. This version is useful for validation but expensive because variant
    selection logic remains on-chain. \\
    Specialized final implementation &
    The optimized implementation produced in this project, where the main protocol modes are
    split into smaller dedicated contracts or clone-based variants. The goal is to move mode
    selection off-chain and avoid paying gas for unused branches. \\
    Hardcoded SHA-256 mode &
    A specialized dispute mode for the case \(\mathsf{desc}=\mathsf{SHA256}\), where the
    circuit is determined by the input length and no generic circuit commitment proof
    \(\pi_1\) is required in Step 8. \\
    \simsox{} / no \(S\)-deposit &
    A simplified optimistic mode in which there is no initial sponsor deposit by \(S\). Each
    participant pays gas in the optimistic part, removing the need for sponsor reimbursement
    logic in that part of the contract. \\
    \bottomrule
  \end{tabularx}
  \caption{Implementation terminology used in the report.}
  \label{tab:intro-terminology}
\end{table}

This distinction is essential for interpreting the experimental results. A number such as
``optimistic cost'', ``deployment cost'', or ``dispute cost'' is meaningful only after saying
which implementation is measured and which protocol steps are included. One of the outcomes
of this project is therefore a stricter accounting discipline: costs are reported either as
one-time infrastructure deployments, per-exchange optimistic costs, dispute initialization
costs, or per-round dispute execution costs.

\section{Project goals}

The project was initially motivated by the need to understand the existing \sox{}
implementation and to make it executable on a complete local stack. After this first phase,
the scope was refined around gas reduction and protocol variants, following discussions with
the supervisor. The final goals can be summarized as follows.

\begin{enumerate}
  \item Reconstruct and validate the execution path of the inherited implementation, from
  off-chain precontract generation to optimistic completion and dispute execution.

  \item Identify which sponsoring variants actually affect gas costs. In particular, the
  project studies no \(S\)-deposit, \(S=B\), \(S=V\), \(S_B=B\), and \(S_V=V\), while
  separating variants determined at precontract time from variants determined only once a
  dispute starts.

  \item Implement the main variants in Solidity without changing the intended protocol
  semantics. The implementation must avoid a ``Swiss-army-knife'' contract where all modes
  are present in one large bytecode and all users pay for unused logic.

  \item Implement and measure a hardcoded \(\mathsf{SHA256}\) dispute mode. For this mode,
  the circuit is fixed by \(|x|\), so the smart contract can reconstruct the relevant gate in
  Step 8 instead of requiring a generic circuit gate and the proof \(\pi_1\).

  \item Produce reproducible gas measurements for the optimistic path and the dispute path,
  including a clear distinction between one-time deployment costs and marginal per-exchange
  costs.

  \item Test large-file execution, especially 1 GiB inputs, using the native Rust command
  line pipeline when the browser/WebAssembly execution path becomes impractical.
\end{enumerate}

These goals are deliberately implementation-oriented. The project does not claim a new
cryptographic proof of the \sox{} protocol. Instead, it asks whether the theoretical variants
that are attractive on paper can be translated into contracts that are measurably cheaper and
still faithful to the intended protocol flow.

\section{Methodology}

The work followed an iterative methodology. The first step was to make the inherited system
run reliably in a local development environment. This required rebuilding the Rust and
WebAssembly components, deploying the Solidity stack, restoring the expected local database
and backend folders, and understanding the actual path used by the frontend. This phase was
necessary because several important costs are not visible by reading a contract alone: the
application decides when a contract is deployed, which role signs which operation, whether a
transaction goes through ERC-4337 tooling, and which off-chain artifacts are stored or
transferred.

The second step was protocol mapping. The implementation uses function names and UI actions
that do not directly match the step names in the paper. For example, the dispute trigger,
challenge response, opinion submission, commitment submission, and finalization functions
correspond to different substeps of the paper protocol. Before measuring gas, these mappings
had to be fixed explicitly. Otherwise, two measurements with the same informal label could
include different pieces of the protocol.

The third step was variant isolation. The initial engineering temptation was to add feature
flags to the existing contracts. This approach is convenient for development because one
contract can simulate many cases. It is not, however, an adequate gas optimization strategy:
Ethereum users pay for bytecode size, storage layout, and runtime branches even when a given
exchange uses only one mode. The project therefore moved progressively from a monolithic
implementation to specialized contracts, clone-based optimistic escrows, and dedicated
dispute games.

The fourth step was measurement. Gas was measured through Hardhat tests designed to execute
complete scenarios rather than isolated functions whenever possible. This matters because
deployment and initialization can dominate the total cost. For large files, off-chain timings
were measured with the native Rust command line tools instead of the browser path, because
the browser/WebAssembly execution path is not the right environment for 1 GiB experiments.
The measurements in the final report distinguish:

\begin{itemize}
  \item one-time infrastructure deployment costs, such as deploying factories and
  implementations;
  \item marginal optimistic costs per exchange, including deployment or clone creation and
  execution until completion;
  \item dispute initialization costs, especially the cost of deploying or creating the
  dispute contract;
  \item post-trigger dispute execution costs, such as challenge rounds and Step 8
  verification;
  \item off-chain timings for precontract generation and dispute witness computation.
\end{itemize}

The final step was consistency checking. Several early measurements were technically correct
but difficult to interpret because they mixed scopes: some came from the monolithic version,
some from the specialized version, some included deployment, and some did not. The report
therefore treats measurement clarity as part of the contribution. A gas table is only useful
if the reader can reconstruct exactly what was paid, by whom, and at which stage of the
protocol.

\section{Main contributions}

The main contributions of this project are the following.

\paragraph{Variant-aware optimistic contracts.}
The optimistic phase was split into dedicated modes for the normal case, no \(S\)-deposit,
\(S=B\), and \(S=V\). In the final specialized implementation, the mode is selected before
deployment by the frontend/backend pipeline, so the contract does not need to carry all
branches at runtime. This is the main reason why the marginal optimistic cost can be reduced
substantially compared with the initial account-per-exchange design.

\paragraph{Self-sponsored dispute execution.}
The dispute contracts were specialized for cases where the dispute sponsors coincide with
the protocol participants, in particular \(S_B=B\) and \(S_V=V\). These cases matter because
they remove unnecessary repetitions of the Step 7--8 loop caused by Step 9 sponsor rotation.
The project therefore separates the cost of one iteration from the cost of two or three
iterations, rather than hiding this effect behind a single aggregate dispute number.

\paragraph{Hardcoded \(\mathsf{SHA256}\) dispute mode.}
For the important case \(\mathsf{desc}=\mathsf{SHA256}\), the circuit is determined by the
input length. A dedicated hardcoded dispute contract was implemented so that Step 8 no
longer requires the generic circuit proof \(\pi_1\). In the strict specialized mode, the
contract reconstructs the relevant SHA-256 gate internally; the vendor does not provide the
generic gate bytes used by the monolithic verifier. This distinction is important because the
monolithic hardcoded variant can remove \(\pi_1\) but still requires a 64-byte gate encoding,
whereas the specialized hardcoded variant removes both.

\paragraph{Large-file execution path.}
The project validated an execution path for large files up to 1 GiB using native command
line tooling. This does not mean that every browser action is suitable for 1 GiB inputs.
Rather, it establishes a practical split: the UI remains useful for demonstration and small
end-to-end exchanges, while the native Rust pipeline is the correct measurement path for
large-file precontract generation and dispute witness computation.

\paragraph{Reproducible gas accounting.}
The project produced benchmark tests for the initial, monolithic, and specialized versions.
It also clarified which costs are one-time protocol deployment costs and which costs are
paid per exchange. This distinction is central to evaluating whether a change is a true gas
optimization or only a relocation of cost from one part of the system to another.

\section{Scope and limitations of the project}

The project focuses on implementation, measurement, and architecture. It does not attempt to
replace the security proof of the \sox{} protocol. It also does not claim that the current
code is production-ready. The implementation remains a research prototype, and several
limitations remain important.

First, the current application still uses a laboratory identity model with local accounts and
development-chain deployment scripts. Wallet integration would be necessary for a realistic
user-facing deployment, but it was not the main objective of this project.

Second, large-file execution is validated through the native command line pipeline, not
through a fully browser-based flow. This is intentional: the original browser/WebAssembly
path is convenient for small demonstrations, but it is not the right tool for measuring 1 GiB
precontract and dispute computations.

Third, the project revealed questions about the inherited circuit encoding and Merkle tree
domain separation. These questions concern the historical circuit format and are not caused
by the gas-specialization work itself, but they are relevant for any future production-grade
version. A clean final design should make gate encodings unambiguous and should separate
Merkle leaves from internal nodes by domain tags.

Finally, some architectural improvements remain future work. In particular, a more radical
factory or registry architecture could avoid deploying a full optimistic account for every
exchange. The present project already moves in that direction with specialized contracts and
clones, but a fully reusable registry-style deployment model would require a broader redesign
of the inherited application.

\section{Report structure}

This first chapter defines the problem, the terminology, the goals, and the methodology of
the project. The remaining chapters will describe the inherited implementation, the protocol
variants, the specialized contract architecture, the experimental measurements, the security
and encoding issues identified during the project, and the installation and demonstration
workflow. The rest of the report will keep the terminology introduced in
Table~\ref{tab:intro-terminology} so that gas costs can be compared without mixing the
initial, monolithic, and specialized implementations.
